\scanpage{10}

\mspage{427}

situation ``limite'', où j'aurais à renoncer à longue échéance à toute activité intellectuelle, qu'elle soit mathématique ou méditante. Plutôt, la tâche pratique la plus immédiate, la plus urgente dans les années à venir, me semble celle justement d'arriver à un équilibre de vie où les deux types d'activité coexistent au jour le jour, celle du corps et celle de l'esprit, sans que l'une ni l'autre ne devienne dévorante et n'évince l'autre. Je ne me cache pas que c'est bien dans la direction ``esprit'' que se trouvent depuis mon enfance mes investissements les plus puissants, que c'est vers elle aussi que me portent aujourd'hui encore les deux principales passions qui ont continué en ces dernières années à dominer ma vie. De ces deux passions, la passion mathématique et la passion de méditation, il me semble que c'est la première nommée surtout, sinon exclusivement, qui agit comme un facteur de déséquilibre dans ma vie - comme quelque chose qui garderait encore une fâcheuse tendance à ``dévorer'' tout le reste au profit d'elle seule. Ce n'est pas un hasard sûrement, si les trois ``\unsure{épisodes maladie|épisodesmaladie}'' dans ma vie qui ont marqué une situation de déséquilibre, depuis juin 1981, se sont placés en des périodes justement où c'est la passion mathématique qui était sur le devant de la scène.

On pourrait dire que ce n'est pas tout à fait le cas pour ce dernier épisode, survenu en cours de rédaction de Récoltes et Semailles, laquelle constitue une période de réflexion sur moi-même, pour ne pas dire une période de méditation à proprement parler. Mais il est vrai aussi que cette réflexion sur mon passé de mathématicien a été alimentée constamment par ma passion mathématique. Il en a été ainsi surtout dans la deuxième partie, l'Enterrement, il me semble, où la composante égotique de cette passion s'est vue impliquée de façon particulièrement forte et constante. Pourtant, même rétrospectivement, je n'ai pas l'impression qu'à aucun moment, cette réflexion ait pris un rythme, un diapason dévorant, voire démentiel, comme en les deux précédentes occasions où mon corps a été contraint finalement à faire entendre un ``ras le bol !'' sans réplique. Vue séparée du contexte de toute \unsure{'une|une} vie, mon activité intellectuelle depuis un an et demi (depuis la ``reprise'' avec la rédaction de la Poursuite des Champs, suivie par Récoltes et Semailles) apparaît comme se poursuivant à un rythme des plus raisonnables, sans y oublier ni le boire ni le manger (mais parfois quand même, un tantinet, le dormir...). Si elle a fini par déboucher sur un troisième ``épisode santé'' (pour utiliser un euphémisme), c'est sans doute sur le fond de toute une vie marquée par ce sempiternel déséquilibre d'une tête trop forte, imposant son rythme et sa loi à un corps robuste qui a longtemps \deleted{\unsure{encaissé} \unsure{sans} \unsure{broncher} (*).}

% archive sheet number: 9

